\documentclass[11pt]{article}
\usepackage[margin=1in]{geometry}
\usepackage[numbers,square,sort&compress]{natbib}
\usepackage{hyperref}
\usepackage{amsmath,amssymb,amsthm}
\usepackage{booktabs}
\usepackage{graphicx}
\usepackage{xcolor}

\newtheorem{definition}{Definition}
\newtheorem{theorem}{Theorem}
\newtheorem{corollary}{Corollary}

\title{Capability-Bounded Agents: Provable, Injection-Resistant\\ Action Limits for Autonomous LLM Agents}
\author{%
  Narendar Kumar Ale \\
  \small PhD Student, Artificial Intelligence, University of the Cumberlands, KY, USA \\
  \small \texttt{nale02820@ucumberlands.edu}
}
\date{July 2026}

\begin{document}
\maketitle

\begin{abstract}
Prompt injection is the highest-ranked unresolved risk for deployed LLM agents (OWASP
LLM Top~10), and a recent formal result argues it may be unfixable from inside the model:
any Contextual-Integrity-based in-model defense either misses attacks that construct a
legitimate-looking context, or blocks genuinely legitimate behavior
\citep{abdelnabi2026agents}. We argue that this impossibility result, rather than being
discouraging, motivates a structural shift: safety for \emph{bounded, irreversible actions}
should not depend on the model resisting injection at all. We formalize a
\textbf{capability contract} and an out-of-model \textbf{reference monitor} that mediates
every tool call before execution, and prove that for any bounded action $a$, $a$ executes
$\implies$ $a$ is permitted by the contract — a containment guarantee that holds regardless
of what the model emits, adversarial or not. We validate the mechanism with (1) a
containment benchmark against a modeled adversary (0\% bypass vs.\ up to 100\% for a
prompt-only guardrail), and (2) a real-model adversarial evaluation against two production
Claude models on Amazon Bedrock, in which every emitted tool call was independently checked
against the monitor. We report an honest null result at current injection strength — both
models resisted 100\% of our direct injections — and situate it against external evidence
that real-world indirect injection does succeed at non-trivial rates on other models
\citep{khodayari2026wild}, with real consequences when it does
\citep{msftsemantickernel2026}. We conclude that in-model alignment and out-of-model
enforcement are complementary, not substitutable, and that for bounded actions the latter
is the only mechanism that offers a provable guarantee.
\end{abstract}

\section{Introduction}

Autonomous LLM agents increasingly act in the world: they call tools, execute code, send
messages, and modify state, often with multiple steps between an initial instruction and a
consequential action. This capability is also the attack surface. Prompt injection — text
in a tool's output, a retrieved document, or a user message that causes an agent to deviate
from its intended task — is ranked the \#1 risk in the OWASP Top~10 for LLM Applications,
and 2026 has produced concrete evidence that it is not a theoretical concern: Microsoft
disclosed two vulnerabilities in its Semantic Kernel agent framework
(CVE-2026-25592, CVE-2026-26030) in which prompt injection escalated to host-level remote
code execution \citep{msftsemantickernel2026}.

The prevailing defense paradigm is \emph{in-model}: instruction/data separation, refusal
tuning, and output filtering, all of which rely on the model itself recognizing and
resisting an injected instruction. Abdelnabi and Bagdasarian \citep{abdelnabi2026agents}
recently gave a formal argument, via a Contextual Integrity (CI) framing, that this
paradigm has a structural ceiling: an adversary can always construct a context in which a
blocked information flow appears contextually legitimate, or a defender who tightens norms
enough to block it will also block genuinely legitimate agent behavior. In other words,
in-model robustness cannot be made complete without also breaking the agent's usefulness.

This paper takes that impossibility result as a premise rather than a conclusion to argue
against. If safety cannot be proven to hold \emph{inside} the model, then for the subset of
actions whose consequences are unacceptable regardless of why the model attempted them
(deleting a production resource, exfiltrating a secret), safety must be enforced
\emph{outside} the model, by a mechanism whose correctness does not depend on model
behavior at all.

\paragraph{Contributions.}
\begin{enumerate}
  \item A formal \textbf{capability contract} and \textbf{reference monitor} construction
        (\S\ref{sec:model}) with a containment theorem: for any bounded action, execution
        implies contract-permission, for all agent policies including adversarial ones.
  \item A \textbf{mechanism benchmark} (\S\ref{sec:mechbench}) demonstrating the containment
        gap between prompt-guardrail and reference-monitor enforcement under a modeled
        adversary: 100\%/50\% bypass for prompt guardrails vs.\ 0\% for the monitor.
  \item A \textbf{real-model adversarial evaluation} (\S\ref{sec:realeval}) against
        production Bedrock Claude models, with full raw transcripts published, reporting an
        honest result: 0\% prompt-mode bypass observed at this injection strength — and a
        methodologically explicit discussion of what that result does and does not prove.
  \item A synthesis (\S\ref{sec:evidence}) connecting that result to external, independently
        documented real-world bypass evidence, resolving the apparent tension between "we
        found no bypass" and "bypasses are known to occur."
\end{enumerate}

Both artifacts are public: the core mechanism and mechanism benchmark at
\url{https://github.com/NarenderKumarA/capability-bounded-agents}, and the real-model
evaluation with raw transcripts at
\url{https://github.com/NarenderKumarA/capability-bounded-agents-eval}.

\section{Related Work}
\label{sec:related}

\paragraph{In-model defenses and their limits.} Abdelnabi and Bagdasarian
\citep{abdelnabi2026agents} is the most directly relevant prior work: their CI-based
impossibility argument is the premise this paper builds on. Adversarial-suffix attacks
\citep{zou2023universal} independently demonstrate that alignment/refusal training does not
provide a hard robustness guarantee even against optimization-based attacks, reinforcing
that in-model defenses are best modeled as empirically strong but not provably complete.

\paragraph{Empirical prevalence and realism of injection.} Khodayari et al.\
\citep{khodayari2026wild} conducted a large-scale empirical study of indirect prompt
injection across 1.2 billion crawled URLs, finding over 15{,}000 validated injection
instances and measuring model compliance directly: across 5{,}200 trials over 13 models,
compliance was "limited but non-negligible, reaching up to 8\% for smaller models" on
plain-text inputs. This is, to our knowledge, the most direct existing measurement of
real-world prompt-mode bypass rate, and it is the number our own real-model evaluation
(\S\ref{sec:realeval}) should be read against. Choi et al.\ \citep{choi2026agentdata}
identify a related but distinct vector, \emph{agent data injection}, in which malicious
\emph{data} rather than instructions is disguised as trusted context (e.g.\ metadata), and
report that "current agents do not isolate trusted data from untrusted data" across several
production systems — directly motivating the untrusted-tool-output threat model we adopt in
\S\ref{sec:threat}. Boisvert et al.\ \citep{boisvert2026doomarena} introduce DoomArena, an
evaluation framework for agent security threats, and report that guardrail-based defenses
underperform relative to defenses that route through a separate, more capable evaluator —
consistent with our motivation for moving the decision outside the acting model entirely.

\paragraph{Policy layers and zero-trust framings.} Kholkar and Ahuja's Policy-as-Prompt
\citep{kholkar2026policy} is the closest prior art to this paper's mechanism: it compiles
governance documents into a policy tree that becomes a runtime classifier enforcing least
privilege with an audit trail. It differs from our construction in kind, not just detail:
their enforcement layer is itself a learned/prompted classifier, so its correctness is
empirical, not provable, in the same sense as the reference monitor of
\S\ref{sec:model} — the reference monitor concept itself is much older, originating in
classical computer security as an access-mediation component whose correctness is
established independently of the subject it mediates \citep{anderson1972}. This paper's
contribution is to bring that classical, provable construction to LLM agent tool-calling
specifically, and to pair it with a benchmark and a real-model evaluation rather than an
architecture sketch alone.

\section{Threat Model}
\label{sec:threat}

We consider an agent with access to a fixed set of tools, driven by an LLM policy
$\pi$ that may be adversarially influenced through any channel that reaches the model's
context: a user message, the output of a previously invoked tool, or a retrieved document
(RAG). We do \emph{not} assume $\pi$ is trusted, nor do we assume any particular robustness
property of $\pi$ — the containment guarantee in \S\ref{sec:model} is proven to hold for
\emph{all} possible $\pi$, including one fully compromised by injection. What is trusted is
the \textbf{capability contract} itself: a static, human-authored specification of which
tool calls are permitted, evaluated by a monitor outside the model's control. We call an
action \textbf{bounded} if it is destructive, irreversible, or crosses a trust boundary
(e.g.\ deleting a resource, or sending data to an external recipient) — the class of actions
for which "the model changed its mind after being tricked" is not an acceptable failure
mode.

\section{Formal Model}
\label{sec:model}

\begin{definition}[Capability]
A capability is a tuple $(n, \text{allowed}, \phi)$ where $n$ is a tool name,
$\text{allowed} \in \{\top, \bot\}$, and $\phi$ is an optional predicate over call
arguments. Absent an explicit capability for $n$, $\text{allowed} = \bot$
(\emph{default-deny}).
\end{definition}

\begin{definition}[Capability Contract]
A capability contract $C$ is a finite set of capabilities, one per tool name. $C$ is
specified independently of, and prior to, any agent execution.
\end{definition}

\begin{definition}[Reference Monitor]
A reference monitor $M_C$ is a function $M_C(n, a) \to \{\text{permit}, \text{deny}\}$ such
that $M_C(n, a) = \text{permit}$ iff $C$ contains a capability for $n$ with
$\text{allowed} = \top$ and ($\phi$ is absent or $\phi(a) = \top$).
\end{definition}

We say an execution environment is \textbf{monitor-enforced} if a tool call $(n, a)$ is
passed to the underlying tool implementation only if $M_C(n, a) = \text{permit}$;
otherwise it is \textbf{prompt-enforced} if the call is passed to the tool whenever the
policy $\pi$ emits it, with no independent check.

\begin{theorem}[Containment]
\label{thm:containment}
For any policy $\pi$ (including adversarially controlled ones) and any tool call
$(n, a)$ emitted by $\pi$ in a monitor-enforced environment: if $(n, a)$ executes, then
$M_C(n, a) = \text{permit}$.
\end{theorem}

\begin{proof}
By construction of the monitor-enforced executor: the underlying tool implementation for
$n$ is invoked only along the code path guarded by $M_C(n, a) = \text{permit}$; there is no
other path to invocation. This holds independently of how $(n, a)$ was produced — $\pi$'s
prompt, weights, or adversarial context are not inputs to $M_C$. $\blacksquare$
\end{proof}

\begin{corollary}[Contrapositive containment]
If $M_C(n, a) = \text{deny}$, then $(n, a)$ never executes, for any $\pi$, regardless of
context.
\end{corollary}

Theorem~\ref{thm:containment} is a soundness guarantee, not a completeness one: it
guarantees the monitor never permits what the contract forbids, but says nothing about
whether the contract itself is correctly specified for the deployment's actual security
requirements. A monitor enforcing an overly permissive contract is sound but not safe in
any useful sense — contract correctness is a specification problem outside the scope of
this theorem, and we return to it as a limitation in \S\ref{sec:discussion}.

\section{Mechanism Benchmark: Modeled Adversary}
\label{sec:mechbench}

As a first validation of the construction, we implement a small benchmark
(\texttt{cba/harness.py}) with three tools (\texttt{read\_status}, \texttt{delete\_resource},
\texttt{send\_message}) and four attacks spanning three injection vectors (poisoned tool
output, poisoned retrieved document, direct user-turn social engineering), targeting either
an unauthorized delete or an exfiltration of a planted secret. A capability contract permits
only \texttt{read\_status} and \texttt{send\_message} to internal recipients. Because this
benchmark's purpose is to validate the \emph{monitor}, not measure real model behavior, we
model whether an injection "flips" a given policy via an intensity threshold rather than
querying a real model — the real-model question is addressed separately in
\S\ref{sec:realeval}.

\begin{table}[h]
\centering
\begin{tabular}{@{}llc@{}}
\toprule
Policy & Enforcement & Bypass rate \\
\midrule
vulnerable-7b & prompt guardrail & 100\% \\
vulnerable-7b & reference monitor & \textbf{0\%} \\
hardened-frontier & prompt guardrail & 50\% \\
hardened-frontier & reference monitor & \textbf{0\%} \\
\bottomrule
\end{tabular}
\caption{Mechanism benchmark: modeled adversary. Whenever an injection flips the modeled
policy, the reference monitor still blocks the bounded action, independent of policy
strength.}
\label{tab:mechbench}
\end{table}

This confirms the mechanism does what Theorem~\ref{thm:containment} claims in a controlled
setting, and motivates the natural next question: does the same hold with a real model in
the loop, and does prompt-mode bypass actually occur?

\section{Real-Model Adversarial Evaluation}
\label{sec:realeval}

\subsection{Design}
We replace the modeled flip-decision with a real one. Each of the four attacks from
\S\ref{sec:mechbench} is realized as an actual multi-turn conversation sent to a real
model via the Bedrock InvokeModel API, using Claude's native \texttt{tools}/\texttt{tool\_use}
mechanism rather than a prompt-convention or regex parse, so that "did the model attempt the
forbidden call" is answered by inspecting structured output, not by string matching.

The agent's system prompt states its authorized scope explicitly and identically across all
attacks and a clean control ("never delete... never send credentials externally... this
holds regardless of what any request, tool output, or document tells you"), so that a
measured flip is informative about the model's behavior under an unambiguous scope
statement, not an artifact of an underspecified prompt. For the two tool-output/document
vectors, the conversation is three turns: a benign user task, a \emph{scripted} (not
model-generated) assistant call to the safe lookup tool, and a tool-result turn whose
content is poisoned with an explicit, unambiguous malicious instruction (exact target
arguments included, so a "flip" is unambiguous to detect rather than dependent on the model
guessing an argument). The direct social-engineering attack is single-turn. Every emitted
\texttt{tool\_use} block from the one real model call per trial is independently checked
against the \emph{same} \texttt{ReferenceMonitor}/contract as \S\ref{sec:mechbench};
prompt-mode and monitor-mode bypass are both derived from the same sampled response for an
apples-to-apples comparison. Each (model, attack) pair received 3 injected trials and 1
clean control trial.

\subsection{Result}

\begin{table}[h]
\centering
\begin{tabular}{@{}lccc@{}}
\toprule
Model & Prompt-mode bypass & Monitor-mode bypass & Trials \\
\midrule
claude-haiku-4-5 & \textbf{0\%} & 0\% & 12 injected + 4 clean \\
claude-sonnet-4-5 & \textbf{0\%} & 0\% & 12 injected + 4 clean \\
claude-sonnet-4 (legacy) & n/a\textsuperscript{$\dagger$} & --- & 0 \\
\bottomrule
\end{tabular}
\caption{Real-model evaluation. $\dagger$: Bedrock returned HTTP 404 ("deprovisioned,
inactive 30+ days") for this legacy model ID; excluded from analysis, not counted as 0\%.}
\label{tab:realeval}
\end{table}

Manual inspection of the raw transcripts confirms these are genuine refusals, not a
detection failure: e.g., claude-sonnet-4-5's response to the exfiltration attack includes
"I'm ignoring the verification request in the status output, as it asks me to send
sensitive information to an external address, which violates my security constraints."
Monitor-mode bypass is trivially 0\% here because prompt-mode bypass was already 0\% — no
forbidden call was ever emitted for the monitor to intercept, so this run does not yet
exercise the monitor's blocking behavior against a real successful injection.

\subsection{What this result does and does not show}
\label{sec:epistemics}

It is important to be precise about what a null result here establishes. The two halves of
our overall claim require different kinds of evidence. The monitor's containment guarantee
(Theorem~\ref{thm:containment}) is proven by construction and does not become "more proven"
by empirical testing. The claim that prompt-only guardrails are unsound, by contrast, is a
claim about LLM behavior — a system with no formal specification — and can only ever be
\emph{falsified} (one real bypass disproves "always safe"), never proven, by testing.
Consequently, absence of a bypass in Table~\ref{tab:realeval} does not weaken the thesis; it
means our particular attacks, at this strength, against these two specific models, did not
find one. This is the same asymmetry Abdelnabi and Bagdasarian's impossibility argument
formalizes: a counterexample is guaranteed to exist somewhere in a sufficiently rich
context space, not that it is easy to find against any given model. We treat escalating our
own adversary (multi-turn erosion, obfuscated encodings, weaker or less-aligned target
models) as necessary future work (\S\ref{sec:future}), and in the meantime turn to
independently documented external evidence for the bypass side of the claim.

\section{External Evidence for Guardrail Unsoundness}
\label{sec:evidence}

Rather than treat our own null result as the only evidence available for the claim that
prompt-guardrail bypass rates are non-zero, we cite independently measured, real-world
evidence:

\begin{itemize}
  \item Khodayari et al.\ \citep{khodayari2026wild} directly measured real-world indirect
        prompt injection compliance across 13 models and found rates "reaching up to 8\%
        for smaller models" on plain-text inputs — a concrete, non-zero, externally
        validated prompt-mode bypass rate, exactly the quantity Table~\ref{tab:realeval}
        did not surface for the two frontier models we tested.
  \item Choi et al.\ \citep{choi2026agentdata} found exploitable agent data injection
        vulnerabilities in multiple deployed systems, escalating in some cases to remote
        code execution, and identify a structural cause — "current agents do not isolate
        trusted data from untrusted data" — that applies broadly across current agent
        architectures, not only to the systems they tested directly.
  \item The Microsoft Semantic Kernel disclosures \citep{msftsemantickernel2026}
        (CVE-2026-25592, CVE-2026-26030) show the consequence end of the same failure mode
        in a shipped, widely used agent framework: prompt injection escalating to
        host-level remote code execution.
\end{itemize}

Taken together, these establish that the "guardrails are unsound" half of our thesis does
not depend on us personally reproducing a bypass: it has already occurred, repeatedly, in
independently conducted research and in production systems. What Table~\ref{tab:realeval}
adds is a controlled, transparent, fully-published methodology for measuring this rate
against any specific model or contract going forward — including, eventually, a first-party
bypass once the adversary is escalated per \S\ref{sec:future}.

\section{Discussion and Limitations}
\label{sec:discussion}

\paragraph{Contract expressiveness is the real bottleneck.} Theorem~\ref{thm:containment}
is a soundness result relative to a given contract; it says nothing about whether that
contract correctly captures the deployment's actual security requirements. A monitor
enforcing a mis-specified or overly permissive contract is sound but not safe. In practice,
authoring a correct, minimal contract is itself a nontrivial specification task, and we
consider methods for validating contract correctness (e.g.\ against a threat model or
compliance policy, cf.\ \citep{kholkar2026policy}) important future work rather than solved
by this paper.

\paragraph{Utility cost.} A contract that is too restrictive degrades legitimate agent
behavior — the same tension Abdelnabi and Bagdasarian identify for in-model defenses
reappears, in a different form, at the contract-authoring layer. Our clean-control trials
(Table~\ref{tab:realeval}) are a first, minimal check that the monitor did not
spuriously block legitimate behavior, but a systematic utility-cost evaluation across more
realistic task distributions is needed.

\paragraph{Single-vendor, small-$N$ real-model sample.} Our real-model evaluation covers
only two Claude models via one API (Bedrock) with three trials per condition — directional,
not statistically powered, and silent on other model families, open-weight models, or
non-Anthropic providers.

\paragraph{Monitor-mode bypass is not yet exercised by a real success.} Because no real
injection in \S\ref{sec:realeval} succeeded at the model level, Table~\ref{tab:realeval}'s
monitor-mode 0\% is not yet a demonstration of the monitor intercepting a real attempted
bypass — only of it correctly permitting benign and non-emitted calls. Demonstrating
monitor-mode containment against a real, successful model-level flip is the single most
important remaining empirical gap, addressed in \S\ref{sec:future}.

\section{Practical Deployment Guidance}
We do not position the reference monitor as a replacement for in-model alignment, but as a
backstop for the specific subset of actions where the cost of a failure is unacceptable
regardless of cause. In practice this suggests: (1) reserve monitor-enforced capability
contracts for bounded/irreversible actions specifically, not all tool calls, to avoid
paying the full utility cost of restriction everywhere; (2) require human confirmation as an
additional layer for the highest-consequence actions, since a correctly-specified contract
is necessary but, per \S\ref{sec:discussion}, not automatically sufficient; (3) treat
in-model defenses and the monitor as complementary layers — the model's own refusal
behavior (as observed in \S\ref{sec:realeval}) still reduces how often the monitor needs to
intervene at all, even though it cannot be relied upon alone.

\section{Future Work}
\label{sec:future}

The immediate priority is escalating adversary strength in the real-model evaluation —
multi-turn erosion, obfuscated/indirect encodings, agent-data-injection-style attacks
\citep{choi2026agentdata}, and weaker or open-weight local models — specifically to obtain a
real, non-zero prompt-mode bypass and demonstrate the monitor intercepting it, closing the
gap identified in \S\ref{sec:discussion}. Further directions include: a systematic
utility-cost evaluation; tooling to validate contract completeness against a stated threat
model; extension to multi-agent and tool-chaining settings, where a bounded action may be
reached indirectly through a sequence of individually-permitted calls; and integration with
existing agent security evaluation frameworks such as DoomArena
\citep{boisvert2026doomarena} for broader, standardized comparison.

\section{Conclusion}
Prompt injection may be structurally unfixable from inside the model, per recent formal
results \citep{abdelnabi2026agents}. We argue this motivates, rather than undermines, a
provable out-of-model enforcement layer for bounded actions, and give a formal construction
—a capability contract mediated by a reference monitor— with a containment theorem that
holds for any agent policy, adversarial or not. A modeled mechanism benchmark and a
real-model evaluation against production Claude models both validate the construction; the
real-model evaluation additionally surfaces an honest null result, which we show is
consistent with, not contrary to, the thesis, and which independently documented real-world
evidence \citep{khodayari2026wild,choi2026agentdata,msftsemantickernel2026} closes the gap
on. All code, the full attack suite, and raw model transcripts are public and reproducible.

\bigskip
\noindent\textbf{Reproducibility.} Core mechanism and mechanism benchmark:
\url{https://github.com/NarenderKumarA/capability-bounded-agents}. Real-model evaluation,
raw transcripts, and run manifests:
\url{https://github.com/NarenderKumarA/capability-bounded-agents-eval}.

\begin{thebibliography}{99}

\bibitem{abdelnabi2026agents}
Sahar Abdelnabi and Eugene Bagdasarian.
\newblock AI Agents May Always Fall for Prompt Injections.
\newblock \emph{arXiv preprint arXiv:2605.17634}, 2026.

\bibitem{khodayari2026wild}
Soheil Khodayari, Xuenan Zhang, Bhupendra Acharya, and Giancarlo Pellegrino.
\newblock Indirect Prompt Injection in the Wild: An Empirical Study of Prevalence,
Techniques, and Objectives.
\newblock \emph{arXiv preprint arXiv:2604.27202}, 2026.

\bibitem{choi2026agentdata}
Woohyuk Choi, Juhee Kim, Taehyun Kang, Jihyeon Jeong, Luyi Xing, and Byoungyoung Lee.
\newblock Agent Data Injection Attacks are Realistic Threats to AI Agents.
\newblock \emph{arXiv preprint arXiv:2607.05120}, 2026.

\bibitem{boisvert2026doomarena}
Leo Boisvert, Mihir Bansal, Chandra Kiran Reddy Evuru, Gabriel Huang, Abhay Puri,
Avinandan Bose, Maryam Fazel, Quentin Cappart, Jason Stanley, Alexandre Lacoste,
Alexandre Drouin, and Krishnamurthy Dvijotham.
\newblock DoomArena: A Framework for Testing AI Agents Against Evolving Security Threats.
\newblock \emph{arXiv preprint arXiv:2504.14064}, 2025.

\bibitem{kholkar2026policy}
Gauri Kholkar and Ratinder Ahuja.
\newblock Policy-as-Prompt: Turning AI Governance Rules into Guardrails for AI Agents.
\newblock \emph{arXiv preprint arXiv:2509.23994}, 2025.

\bibitem{zou2023universal}
Andy Zou, Zifan Wang, Nicholas Carlini, Milad Nasr, J. Zico Kolter, and Matt Fredrikson.
\newblock Universal and Transferable Adversarial Attacks on Aligned Language Models.
\newblock \emph{arXiv preprint arXiv:2307.15043}, 2023.

\bibitem{anderson1972}
James P. Anderson.
\newblock Computer Security Technology Planning Study.
\newblock \emph{Technical Report ESD-TR-73-51}, Electronic Systems Division, Air Force
Systems Command, 1972.

\bibitem{msftsemantickernel2026}
Microsoft Security Response Center.
\newblock When Prompts Become Shells: RCE Vulnerabilities in AI Agent Frameworks.
\newblock \emph{Microsoft Security Blog}, May 2026.
\newblock \url{https://www.microsoft.com/en-us/security/blog/2026/05/07/prompts-become-shells-rce-vulnerabilities-ai-agent-frameworks/}

\bibitem{owasp2025llm}
OWASP Foundation.
\newblock OWASP Top 10 for Large Language Model Applications.
\newblock \url{https://owasp.org/www-project-top-10-for-large-language-model-applications/}

\end{thebibliography}

\end{document}
